% brief.tex — short self-contained reference (target: 15–30 pages)
%
% USAGE BY CLAUDE:
%   - READ THIS FIRST at the start of every new session.
%   - This file contains the current state of the project: established results,
%     key formulas, open problems, conventions. No proofs, no derivations.
%   - Cross-references to workbook.tex are provided so you can navigate there when
%     you need the derivation.
%   - Update this file whenever workbook.tex gets a new result, corrected formula,
%     or status change. Use /sync-condensed to propagate load-bearing changes.
%   - Status tags: [ESTABLISHED], [CONJECTURED], [OPEN]
%
\documentclass[11pt, a4paper]{article}

\usepackage{amsmath, amssymb, amsthm}
\usepackage{mathtools}
\usepackage[margin=2.5cm]{geometry}
\usepackage{microtype}
\usepackage[T1]{fontenc}
\usepackage[utf8]{inputenc}
\usepackage[colorlinks=true, linkcolor=blue, citecolor=blue, urlcolor=blue]{hyperref}
\usepackage{enumitem}
\usepackage{xcolor}
\usepackage{mdframed}           % shaded boxes for highlighted results

% Status colours
\newcommand{\established}{\textcolor{teal}{\textbf{[ESTABLISHED]}}}
\newcommand{\conjectured}{\textcolor{orange}{\textbf{[CONJECTURED]}}}
\newcommand{\open}{\textcolor{red}{\textbf{[OPEN]}}}

\theoremstyle{plain}
\newtheorem{theorem}{Theorem}[section]
\newtheorem{proposition}[theorem]{Proposition}
\newtheorem{lemma}[theorem]{Lemma}
\newtheorem{conjecture}[theorem]{Conjecture}

\theoremstyle{definition}
\newtheorem{definition}[theorem]{Definition}
\newtheorem{remark}[theorem]{Remark}

% -----------------------------------------------------------------------
% CUSTOM COMMANDS — keep in sync with workbook.tex and CLAUDE.md
% -----------------------------------------------------------------------
% \newcommand{\examplecommand}[1]{\mathcal{#1}}

% -----------------------------------------------------------------------
\title{[Project Title] — Condensed Reference\\
{\normalsize\textit{Current state of knowledge — updated \today}}}
\author{[Author Name]}
\date{}
% -----------------------------------------------------------------------

\begin{document}

\maketitle

\tableofcontents
\newpage

% -----------------------------------------------------------------------
\section{Main results}
\label{sec:mainresults}
% -----------------------------------------------------------------------

% Two or three key theorems stated precisely.
% Someone reading only this section should know what the project has established.

\begin{theorem}[\established, see workbook.tex \S\ref{...}]
[Statement of main result.]
\end{theorem}

% -----------------------------------------------------------------------
\section{Conventions and definitions}
\label{sec:conventions}
% -----------------------------------------------------------------------

% Every symbol defined, with the exact normalisation used.
% This section prevents Claude from guessing. Be explicit about signs,
% normalisations, and branch choices.

\begin{definition}
[Define your key objects here.]
\end{definition}

\begin{remark}
[Normalisation convention: state explicitly which convention you use when
there are multiple in the literature.]
\end{remark}

% -----------------------------------------------------------------------
\section{[First topic] — results}
\label{sec:topic1results}
% -----------------------------------------------------------------------

% One section per topic. Theorem statement only, no proof.
% Cross-reference workbook.tex for the derivation.
% Tag each result: [ESTABLISHED], [CONJECTURED], or [OPEN].

\begin{proposition}[\established, workbook.tex \S\ref{...}]
[Statement.]
\end{proposition}

\begin{conjecture}[\conjectured]
[Statement.]
\end{conjecture}

% -----------------------------------------------------------------------
\section{[Second topic] — results}
\label{sec:topic2results}
% -----------------------------------------------------------------------

[Content following same pattern.]

% -----------------------------------------------------------------------
\section{Numerical results}
\label{sec:numericresults}
% -----------------------------------------------------------------------

% Exact computed values, method, precision, what was validated.
% For each result: value, method, precision level, validation performed.

\begin{center}
\begin{tabular}{llll}
\toprule
Quantity & Value & Method & Validated \\
\midrule
[quantity] & [value] & [method] & [how validated] \\
\bottomrule
\end{tabular}
\end{center}

% -----------------------------------------------------------------------
\section{Open problems}
\label{sec:open}
% -----------------------------------------------------------------------

% Ranked list. Item 1 is the next thing to do.

\begin{enumerate}[label=\textbf{O\arabic*.}]
  \item \open\ [Most urgent open problem — this is the next step.]
  \item \open\ [Second priority.]
  \item \open\ [Third priority.]
\end{enumerate}

% -----------------------------------------------------------------------
\section{Recent corrections}
\label{sec:corrections}
% -----------------------------------------------------------------------

% Any result that was wrong and has been corrected.
% State the corrected version; note what the previous version was.

% [No corrections yet.]

% -----------------------------------------------------------------------
\end{document}
